\section{Sterowalność}

\subsection{Macierz sterowalności}

\begin{equation}
\mathcal{C} = \begin{bmatrix} B & AB & A^2 B \end{bmatrix}
\end{equation}

Macierz sterowalności przyjmuje postać:

\begin{equation}
\mathcal{C} = \begin{bmatrix} 0 & 0 & 12{,}389 \\ 0 & 12{,}389 & -110{,}796 \\ 10{,}526 & -110{,}796 & 1{,}166 \times 10^3 \end{bmatrix}
\end{equation}

\subsection{Rząd macierzy sterowalności}

Aby zbadać rząd macierzy, oblicza się jej wyznacznik:

\begin{equation}
\det(\mathcal{C}) = 0 \cdot \begin{vmatrix} 12{,}389 & -110{,}796 \\ -110{,}796 & 1{,}166 \times 10^3 \end{vmatrix} - 0 \cdot \begin{vmatrix} 0 & -110{,}796 \\ 10{,}526 & 1{,}166 \times 10^3 \end{vmatrix} + 12{,}389 \cdot \begin{vmatrix} 0 & 12{,}389 \\ 10{,}526 & -110{,}796 \end{vmatrix}
\end{equation}

Obliczając podwyznacznik:

\begin{equation}
\begin{vmatrix} 0 & 12{,}389 \\ 10{,}526 & -110{,}796 \end{vmatrix} = 0 \cdot (-110{,}796) - 12{,}389 \cdot 10{,}526 = -130{,}380
\end{equation}

Stąd:

\begin{equation}
\det(\mathcal{C}) = 12{,}389 \times (-130{,}380) \approx -1614
\end{equation}

Ponieważ $\det(\mathcal{C}) \neq 0$, macierz sterowalności ma pełny rząd 3. Układ jest w pełni sterowalny co jest warunkiem koniecznym do zaprojektowania sprawnego regulatora.
